%! Function Fundamentals — Unit 1, Lesson 1.1 — Exit Ticket (Answer Key)
\documentclass[10pt]{article}
\usepackage{precalculus-article}
\usepackage{precalculus-key}   % pulls in -boxes; do NOT also load -boxes
\newcommand{\TallMath}[1]{$\displaystyle #1\rule[-1.4em]{0pt}{3.2em}$}

\begin{document}

\pageheader{Unit 1, Lesson 1.1}{Exit Ticket --- Answer Key}

\vspace{0.12in}

\begin{enumerate}[itemsep=12pt]

\item Let $f(x) = 4x - 3$.

\begin{work}
  f(2) &= 4(2) - 3 \\
       &= 5
\end{work}

$f(2) = $ \ans{5} \qquad $f(0) = $ \ans{$-3$}

\item The table shows a function $g$.

\smallskip
\begin{tabular}{|c|c|c|c|c|}
\hline
$x$ & 0 & 1 & 2 & 3 \\
\hline
$g(x)$ & 7 & 10 & 13 & 16 \\
\hline
\end{tabular}
\smallskip

\begin{enumerate}[label=(\alph*), itemsep=4pt]
  \item $g(2) = $ \ans{13}
  \item Solve $g(x) = 10$: \quad $x = $ \ans{1}
\end{enumerate}

\item The function $f = \{(1, 4),\ (2, 7),\ (3, 7),\ (5, 9)\}$.

\begin{enumerate}[label=(\alph*), itemsep=4pt]
  \item Domain: \ans{$\{1, 2, 3, 5\}$}
  \item Range: \ans{$\{4, 7, 9\}$ (the repeated 7 is listed once)}
\end{enumerate}

\end{enumerate}

\end{document}
